% Native semantic reconstruction of source physical page 35 (printed page 34).
% Provenance: IEEE Std 1516-2025, Federate Rules 9 and 10.
\HLAStartSemanticPage{35}{34}
\fontsize{10.02pt}{11.566pt}\selectfont
\begingroup
\hbadness=10000\relax
\addfontfeatures{FakeStretch=1.035}
\vspace*{-0.85pt}

\HLADetailedRuleHeading{detailed-rule-9}{6.4 Rule 9}
\vspace{12.77pt}

\HLANormativeRuleStatement{Federates shall be able to vary the conditions (e.g., thresholds) under which they provide updates of\linebreak instance attributes, as specified in their SOMs.}
\vspace{11.87pt}

The HLA permits federates to own (i.e., provides the privilege to produce updated values for) instance\linebreak
attributes of object instances represented in the federate and to then make those values available to other\linebreak
federates through the RTI. Different federations may specify different conditions under which instance\linebreak
attributes will be updated [e.g., at some specified rate, or when the amount of change in value exceeds a\linebreak
specified threshold (such as altitude changes of more than 1000 ft)]. The conditions applicable to the update\linebreak
of specific instance attributes owned by a federate shall be documented in the SOM for that federate.\par
\vspace{28.90pt}

\HLADetailedRuleHeading{detailed-rule-10}{6.5 Rule 10}
\vspace{12.77pt}

\HLANormativeRuleStatement{Federates shall be able to manage local time in a way that will allow them to coordinate data exchange\linebreak with other members of a federation.}
\vspace{11.87pt}

Federation designers will identify their time management approach as part of their implementation design.\linebreak
Federates shall adhere to the time management approach of the federation.\par
\endgroup
\clearpage
